% ============================================================
%  §1  Introduction
% ============================================================

\section{Introduction}
\label{sec:introduction}

% ---- scope declaration ----
This paper operates within the VED$\times$IFGT framework
\citep{VED_vol1, IFGT} and applies it to the domain of
conceptual structure and thought dynamics.
All variables follow the Master Variable Table v2.1.
The framework developed here is called Conceptual Topography (CT);
the phrase ``conceptual terrain'' refers to the same object when the
intuitive metaphor is emphasized.

% ---- motivation ----
\subsection{The Problem of Conceptual Structure}

What is a concept?
The standard answer places concepts in the semantic domain:
a concept is a meaning-bearing unit, a mental representation,
or a node in a propositional network.
This answer, while useful for many purposes, leaves the
\emph{dynamical} question unanswered.
How does a concept form, stabilize, and constrain subsequent thought?
How does understanding arise, and why does it sometimes fail to transfer
even when the relevant definitions are available?

The present paper approaches these questions from a different direction.
Rather than asking what a concept \emph{means}, we ask what structural
conditions allow a region of the information field to persist, be
re-entered at low cost, and systematically bias future dynamics.
This reframing shifts the problem from semantics to geometry.

% ---- position ----
\subsection{Position in the \texorpdfstring{VED$\times$IFGT}{VED x IFGT} Series}

The VED$\times$IFGT framework consists of two foundational layers.
Vortical Enclosure Dynamics (VED) describes the generation of structure
from difference, gradient, flow, and closure \citep{VED_vol1}.
Information Field Geometry Theory (IFGT) describes the complementary
informational structure that persists as quasi-closure in the region
where $C$ does not reach unity \citep{IFGT}.

The present work is positioned as \emph{Application Layer 1} of this
framework: it applies the field-theoretic structure of IFGT to the
emergence of conceptual landscapes in cognitive and social systems.
It does not introduce new foundational variables; it identifies the
solution classes of the unified VED$\times$IFGT equation that
correspond to cognitive phenomena.

The series architecture is:
\begin{center}
\begin{tabular}{ll}
\toprule
Layer & Theory \\
\midrule
L0 (generation) & VED \\
L0 (structure)  & IFGT \\
L1 (application) & \textbf{Conceptual Topography [this paper]} \\
L1 (application) & Consciousness Theory [forthcoming] \\
L1 (application) & Intelligence Theory [forthcoming] \\
\bottomrule
\end{tabular}
\end{center}

% ---- main claims ----
\subsection{Main Claims}

The paper establishes the following:

\begin{enumerate}
  \item A \textbf{concept} is a locally stable region of the information
        potential $\Phi(x,\tau)$, defined by $\nabla\Phi(x)\approx 0$
        with stability under small perturbations.

  \item \textbf{Understanding} is a temporary reduction in traversal cost
        across a conceptual region, not the acquisition of correct meaning.

  \item \textbf{Thought} is the trajectory of information flow $\JI$
        across the potential landscape; it is constrained by $\nabla\Phi$
        and does not require a directing subject.

  \item \textbf{Language tokens} function as discrete re-addressable units
        that couple locally to $\nabla\Phi$, deepening attractors through
        repeated use without carrying semantic content.

  \item The strict positivity $I(x,\tau)>0$ excludes complete
        conceptual closure as an ordinary dynamical state and leaves the
        landscape open to further differentiation.

  \item At the \textbf{social scale}, the same dynamics produce a shared
        potential geometry maintained across agents, constituting the
        structural basis of norms, institutions, and collective intelligence.
\end{enumerate}

% ---- structure of the paper ----
\subsection{Structure of the Paper}

Section~\ref{sec:intuitive} presents an intuitive picture of conceptual
landscapes.
Section~\ref{sec:formal} establishes the formal mapping from VED$\times$IFGT
variables to conceptual phenomena.
Section~\ref{sec:core} derives the core equations governing landscape
formation and token--field coupling.
Section~\ref{sec:structural} provides structural interpretation of the
main solution classes.
Section~\ref{sec:dynamics} analyzes the dynamics of sedimentation,
forgetting, and conceptual change.
Section~\ref{sec:multiscale} treats the scale inversion from individual
to social landscapes.
Section~\ref{sec:pathologies} identifies failure modes and limits.
Section~\ref{sec:note} connects the formal structure to the intuitive
descriptions developed in the companion note series.
Section~\ref{sec:discussion} discusses implications and open questions.
Section~\ref{sec:conclusion} concludes.
